\section{Introduction}
\label{sec:intro}

LLM agents increasingly act for people where other agents act for other people: they schedule, procure and edit shared code \citep{cooperbench,parallact}. When the principals want different things, someone has decided how the agent handles the conflict, usually in a system prompt (``be firm'', ``be collaborative'') and usually by intuition. This paper measures what that choice is worth.

\begin{figure}[t]
\centering
\FigOverview
\caption{(a) The five conflict-handling strategies of the \tki{} model. (b) What each strategy, given as a rule, does to the default outcome when one agent adopts it against a default (\neutral) partner in the core regime (E1, 80 games per arrow): each arrow runs from the default pair's payoffs to the adopter's and its partner's. Above the dashed line the joint payoff rises.}
\label{fig:overview}
\end{figure}

Organizational psychology has a vocabulary for the choice. The \tki{} model \citep{thomas1974tki,thomas1992conflict} places conflict behaviour on two dimensions, assertiveness and cooperativeness, and names five strategies (\cref{fig:overview}a). It also says when each fits \citep{thomas1977toward}: compete when quick, decisive action is vital; collaborate when both sides' concerns are too important to be compromised; compromise under time pressure or when goals are mutually exclusive; accommodate when the issue matters more to the other side; avoid when confrontation costs more than it yields. These are claims about outcomes. In people the styles are measured by self-report; in LLM agents a strategy can be assigned, its enactment checked and its consequences measured exactly.

We do this in three steps. First, a \emph{parametric multi-issue bargaining environment}: two agents with private weights settle a design over several issues by alternating offers, the Pareto frontier is computed exactly, and one knob each sets urgency, the cost of time and of talk, stakes, power, information and how much there is to trade (\cref{sec:env}). Second, each strategy as \emph{behavioural rules} on the agent's own point scale, so that compliance is checked from the moves without a judge. Third, about 20{,}000 games: the full strategy matrix, ten situational conditions, three realistic framings of the same instances, 85 instances derived from real conflicts between code changes, two published negotiation benchmarks, four models and pairs of different models (\cref{sec:design}). We propose no method; the contribution is the map.

\paragraph{Findings.} \emph{(i)} Rules are enacted and style sentences mostly are not (\cref{sec:r-enact}). \emph{(ii)} In the core regime, \collaborating{} adopted alone earns its user 8.4 points of 100 more than the default agent earns while the partner keeps what it had (\cref{fig:overview}b); it is the best response to four of six partners and the attractor of replicator dynamics. \competing{} gains nothing significant and costs the partner 11 points, two accommodators are the most efficient pair, and \compromising{} and \avoiding{} destroy value (\cref{sec:r-matrix}). \emph{(iii)} The situation moves the map, but not as the textbook prescribes: of the six conditions for which the model names a strategy, that strategy is the best one to adopt in one (\cref{sec:r-map}). \emph{(iv)} Cost orders the strategies the other way round from payoff, and a modest price on tokens hands the answer to the cheap strategies (\cref{sec:r-cost}). \emph{(v)} Across 29 settings the negative claims transfer and the positive one is conditional: a \collaborating{} rule raises the joint payoff with almost every partner, and who keeps the gain is decided by the partner (\cref{sec:r-transfer}).

\paragraph{Contributions.} (1) A bargaining environment for LLM agents with exact frontiers and knobs for the situational variables the \tki{} literature names, in three layers: a synthetic core, scenario skins, and instances derived from real code conflicts. (2) A rule-based operationalization of the five strategies with a compliance checker, compared with prompted style sentences. (3) A payoff map (the strategy matrix, the condition map and the price of tokens at which each answer changes) and a transfer analysis over framings, real-conflict structures, two external benchmarks, four models and cross-model pairs, with every configuration and its predictions frozen before the run.

